%% kapa: exam=1 round=33 year=2022 period=2 form=A subject=accounting subject_ko=회계학 seq=17 number=57
%% source: Q-Net 감정평가사 기출문제 (cst00309) artlSeq=5209997 file=raw/1008/5209997/2022년도 제33회 감정평가사 1차시험 2교시(감정평가 관계법규, 회계학).pdf
%% conversion: kordoc 4.24.51 -> tools/gen1.py (auto); review=proofread-sonnet; flags=none
\begin{question}{57}
금융상품에 관한 설명으로 옳지 않은 것은?
\begin{choices}
  \item 금융자산의 정형화된 매입 또는 매도는 매매일이나 결제일에 인식하거나 제거한다.
  \item 당기손익-공정가치 측정 금융자산이 아닌 경우 해당 금융자산의 취득과 직접 관련되는 거래원가는 최초 인식시점의 공정가치에 가산한다.
  \item 금융자산의 계약상 현금흐름이 재협상되거나 변경되었으나 그 금융자산이 제거되지 아니하는 경우에는 해당 금융자산의 총 장부금액을 재계산하고 변경손익을 당기손익으로 인식한다.
  \item 금융자산 양도의 결과로 금융자산 전체를 제거하는 경우에는 금융자산의 장부금액과 수취한 대가의 차액을 당기손익으로 인식한다.
  \item 최초 발생시점이나 매입할 때 신용이 손상되어 있는 상각후원가 측정 금융자산의 이자수익은 최초 인식시점부터 총 장부금액에 유효이자율을 적용하여 계산한다.
\end{choices}
\end{question}
\answer{5}
